\documentclass[10pt]{article}

% Compilar con LuaLaTeX o XeLaTeX
\usepackage{fontspec}
\setmainfont{TeX Gyre Pagella} % si no está, cámbiala
\setsansfont{Fira Sans}       % instalar o cambiar
\setmonofont{DejaVu Sans Mono}       % instalar o cambiar

\usepackage[spanish]{babel}
\usepackage{microtype}
\usepackage{geometry}
\geometry{margin=2.5cm}
\usepackage{xcolor}
\definecolor{unabblue}{HTML}{003366}
\definecolor{unabgold}{HTML}{C69214}
\definecolor{lightblue}{HTML}{E6F0FF}

\usepackage[most]{tcolorbox}
\tcbset{enhanced,
  colback=lightblue!30,
  colframe=unabblue,
  fonttitle=\sffamily\bfseries\large,
  boxrule=1pt,arc=3mm,
  left=2mm,right=2mm,top=1mm,bottom=1mm,
  before skip=8pt,after skip=12pt
}

\usepackage{fancyhdr}
\pagestyle{fancy}
\fancyhf{}
\lhead{\sffamily\bfseries\textcolor{unabblue}{Prueba de Programación -- Versión A}}
\rhead{\sffamily\textcolor{unabblue}{PCFI161}}
\cfoot{\thepage}

\usepackage{amsmath,amssymb}
\usepackage{hyperref}
\hypersetup{colorlinks=true,linkcolor=unabblue,urlcolor=unabblue}

\begin{document}
\begin{center}
  {\huge\bfseries\textcolor{unabblue}{Prueba de Programación}}\\[4pt]
  {\Large\sffamily\textcolor{unabgold}{Física y Astronomía}}\\[6pt]
  {\large\sffamily Versión A}\\[10pt]
  \today
\end{center}

\vspace{0.5em}
\noindent\textbf{Instrucciones}:

\begin{itemize}
    \item Debe contestar en total \textbf{sólo 3 preguntas}.
    \begin{itemize}
        \item \textbf{Una} de ellas debe ser elegida entre las dos primeras (mayor dificultad).
        \item \textbf{Dos} de ellas deben ser elegidas entre las 3 últimas (menor dificultad)
    \end{itemize}
    \item Si responde más de las preguntas solicitadas, sus respuestas se anularán y se le asignará la nota 1.0.
\end{itemize}

\noindent\makebox[\linewidth]{\rule{\textwidth}{0.4pt}}
\noindent\makebox[\linewidth]{\rule{\textwidth}{0.4pt}}
\section*{Preguntas con mayor dificultad}
% -------------------- Pregunta 1 ---------------------------
\begin{tcolorbox}[title={Pregunta 1 – Órbita elíptica parametrizada}]
  En una elipse el radio vector en función del ángulo $\theta$ es  
\[
r(\theta) = \frac{a\cdot (1-e^2)}{1+e\cdot \cos(\theta)}
\]
  donde  
  \begin{itemize}
      \item $a = $1.0 UA  (UA = unidad astronómica, dato fijo)
      \item $e= $0.40 (número adimensional llamado \textit{excentricidad}, si $e=0$ es un círculo)
  \end{itemize}

   Se pide:
   \begin{itemize}
       \item Crear un \textit{array} (arreglo) \texttt{theta} con 720 valores uniformes en $[0, 2\pi)$ mediante \texttt{np.linspace}.
       \item Calcular \texttt{r} con la ecuación $r(\theta)$ y luego obtener $x = r\cdot \cos(\theta)$, y $y=r\cdot\sin(\theta)$, usando NumPy (sin bucles).
       \item Graficar todos los puntos $(x, y)$ en una ventana cuadrada (\texttt{plt.axis('equal')})
       \item Con NumPy verificar que los promedios $\langle x \rangle$ e $\langle y \rangle$ son $\sim 0$.
   \end{itemize}
\end{tcolorbox}

% -------------------- Pregunta 2 ---------------------------
\begin{tcolorbox}[title={Pregunta 2 – Mini-catálogo HIPPARCOS}]
Se dispone del CSV:\\

{\footnotesize
\url{https://gitarra.cl/lectures/gfiles/-/raw/main/pcfi161/S12/stars_brightness.csv}}\\


Columnas: \texttt{spectral\_class}, \texttt{temperature\_K}(temperatura [K]), \texttt{magnitude\_app}(mag. aparente).\\

Se pide:
\begin{itemize}
    \item Leer el archivo en un \texttt{DataFrame}
    \item Para cada \texttt{spectral\_class} calcular y mostrar
    \begin{enumerate}
        \item Temperatura media $\langle$\texttt{temperature\_K}$\rangle$.
        \item Desviación estándard de \texttt{temperature\_K}.
        \item Magnitud aparente mínima (estrella más brillante)
    \end{enumerate}
    \item Ordenar la tabla por $\langle$\texttt{temperature\_K}$\rangle$ descendente
    \item Graficar \texttt{temperature\_K} \textit{vs} \texttt{magnitude\_app} (scatter); comentar brevemente la relación que se observa (¿estrellas más calientes son más brillantes?)
\end{itemize}

\end{tcolorbox}

\noindent\makebox[\linewidth]{\rule{\textwidth}{0.4pt}}
\noindent\makebox[\linewidth]{\rule{\textwidth}{0.4pt}}
\section*{Preguntas con menor dificultad}
% -------------------- Pregunta 3 ---------------------------
\begin{tcolorbox}[title={Pregunta 3 – Verificación directa de la ley de Kepler}]
Datos (planetas interiores del Sistema Solar):
\begin{itemize}
    \item Semieje mayor: \texttt{a} [UA]      = [0.39, 0.72, 1.00, 1.52]
    \item Periodo orbital: \texttt{T} [años]  = [0.24, 0.61, 1.00, 1.88]
\end{itemize}

Para una órbita casi circular la teoría predice:
\[
      T = k \cdot a^{(3/2)}
\]

donde $k$ debería ser 1 si $T$ se mide en años y $a$ en UA.

   Se pide:
\begin{itemize}
    \item Crear un gráfico tipo scatter con $a$ en el eje $x$ y $T$ en el eje $y$.
    \item Para cada planeta calcule el cociente $R_i = \frac{T_i}{(a_i)^{(3/2)}}$. Muestre los cuatro valores de $R_i$ con 3 decimales.
    \item Calcule el promedio $R_{prom}$ y la desviación estándar $\sigma_R$ de los $R_i$.
\end{itemize}
\end{tcolorbox}

% -------------------- Pregunta 4 ---------------------------
\begin{tcolorbox}[title={Pregunta 4 – Filtrado rápido de exoplanetas}]
  Se dispone de un CSV :\\


{\footnotesize
\url{https://gitarra.cl/lectures/gfiles/-/raw/main/pcfi161/S13/exoplanets_full.csv}}\\
   
con columnas: \texttt{name}, \texttt{mass\_Mj} (masas de Júpiter), \texttt{period\_d} (días).\\

Se pide:
\begin{itemize}
    \item Cargar el archivo en un \texttt{DataFrame}.
    \item Filtrar filas con \texttt{mass\_Mj}$\leq$1.5, y \texttt{period\_d}$\leq$10.
    \item Genere un resumen estadístico de la tabla, por ejemplo promedios, desviaciones, etc.
\end{itemize}
\end{tcolorbox}

% -------------------- Pregunta 5 ---------------------------
\begin{tcolorbox}[title={Pregunta 5 – Razón de Fibonacci y número áureo}]
Se pide:
\begin{itemize}
    \item Generar los primeros 20 términos de la serie de Fibonacci:
    $$F_1=1\qquad F_2=1\qquad F_n = F_{n-1} + F_{n-2}$$
    usando sólo dos variables y un buvle \texttt{for}.
    \item Para $n \geq 2$ imprimir la razón $F_n/F_{n-1}$ con 6 decimales.
    \item Mostrar el valor final calculado y compararlo con el valor teórico $\phi\sim 1.618034$.
\end{itemize}
\end{tcolorbox}


\end{document}

